\documentclass[11pt]{article}
\usepackage{amsmath,amssymb,amsthm}
\newtheorem{theorem}{Theorem}
\newtheorem{lemma}{Lemma}
\title{Two compact reverse heads shatter five columns}
\author{}
\date{}
\begin{document}
\maketitle

\paragraph{Status and scope.}
This is a complete written construction, with an independent exact-rational
checker. The fixed five-query compiler and its strict axis-zero consequence
are now Lean-checked in \texttt{Results/ReverseFiveQueries.lean}. The
arbitrarily-small-slope strengthening below remains written only. It gives a fixed constant improvement in a
local grid representation, not either accepted asymptotic target. In
particular, no amplification to $\Omega(H\log H)$ shattered columns is
proved here. All denominators below are positive on the entire Boolean
cube and have strictly negative slopes. The first theorem does not impose
axis zeros; the existing checked axis-repair theorem supplies that step.

\begin{theorem}
For every positive integer $r$ and every sign table
$\sigma:[r]\times[5]\to\{-1,1\}$, there are two affine ratios
whose sum has strict sign $\sigma(i,j)$ at $(e_i,e_j)$. Their denominators
are strictly reverse and positive on the full Boolean $(r+5)$-cube.
The sum of the absolute denominator slopes can be made arbitrarily
small while the denominator intercepts stay equal to one.
\end{theorem}

\begin{proof}
Fix $0<\epsilon\le 1/100$, put $t=1+\epsilon$, and first work with
the five normalized queries
\[
\begin{array}{c|cc|cc}
j&D_{1j}&D_{2j}&b_{1j}&b_{2j}\\ \hline
0&1&1&0&0\\
1&t&t&0&0\\
2&1-\epsilon&1&\epsilon(1-\epsilon)&-\epsilon\\
3&1-\epsilon&1&-\epsilon(1-\epsilon)&\epsilon\\
4&1-2\epsilon&1&4\epsilon(1-2\epsilon)&-4\epsilon.
\end{array}
\]
The score at query $j$ is
\[
 f_j=\frac{a_1+b_{1j}}{D_{1j}-p_1}
       +\frac{a_2+b_{2j}}{D_{2j}-p_2}.
\]
Consider one desired sign vector $s_0,\ldots,s_4$. If $s_0=1$, set
$Y=(8s_2,s_3,s_4)$; if $s_0=-1$, set $Y=(s_2,8s_3,s_4)$.
Define
\[
 X=\frac{Y_1+Y_2}{2},\qquad
 v=\frac{5Y_1-3Y_2-2Y_3}{4},\qquad
 P=\frac{2Y_3-3Y_1+Y_2}{4\epsilon},\qquad Q=Xv.
\]
Thus $X+v+\epsilon P=Y_1$,
$X-v-\epsilon P=Y_2$, and $X+2v+4\epsilon P=Y_3$.
Moreover $Q$ has sign $s_0$: when $|Y_1|=8$, both factors
in $8Q=(Y_1+Y_2)(5Y_1-3Y_2-2Y_3)$ have sign $Y_1$;
when $|Y_2|=8$, their signs are opposite. The estimates
\[
 |X|\le5,\quad |v|\le12,\quad |P|\le7/\epsilon,
 \quad |P+v|\le8/\epsilon,\quad 0<|Q|\le60
\]
follow directly. Let
\[
 \lambda=\begin{cases}0&s_0=s_1,\\
 -Q(t+1)/(2t)&s_0\ne s_1.
 \end{cases}
\]
For a sufficiently small positive $\delta$, set
\[
 p_2=\delta P,\quad p_1=\delta(P+v),\quad
 a_1=\delta X,\quad a_2=-\delta X+\delta^2\lambda.
\]
Clearing the two positive denominators gives, exactly,
\[
\begin{split}
 N_0&=\delta^2(\lambda+Q)-\delta^3\lambda(P+v),\\
 N_1&=\delta^2(t\lambda+Q)-\delta^3\lambda(P+v).
\end{split}
\]
If $s_0=s_1$, these have the required sign because $\lambda=0$.
Otherwise their leading coefficients are
$Q(t-1)/(2t)$ and $-Q(t-1)/2$. Both signs are preserved whenever
$8\delta/\epsilon<\epsilon/(2t)$, since $|\lambda|\le|Q|$.
For queries $j=2,3,4$, write respectively
$k_j=\epsilon,\epsilon,2\epsilon$. Their cleared numerators are
\[
 N_j=\delta k_jY_{j-1}
   +\delta^2\bigl((1-k_j)\lambda+Q\bigr)
   -\delta^3\lambda(P+v).
\]
Thus the desired remaining signs hold, for example, if
$200\delta<\epsilon$ and $8\delta/\epsilon<1$.

Here is an explicit embedding that satisfies all these inequalities
simultaneously for all rows and the entire Boolean cube. Put
\[
 \rho=\frac{\epsilon}{100r},\quad
 \delta=\frac{\rho\epsilon^2}{100},\quad d_0=1-10\epsilon.
\]
Choose the above parameters independently in every row. Define
\[
 \widehat p_{1i}=\rho+d_0\delta(P_i+v_i),\quad
 \widehat p_{2i}=\rho+d_0\delta P_i,\qquad
 \widehat q_{hj}=1-\rho-d_0D_{hj}.
\]
The actual denominators and numerators are
\[
 B_h(z,w)=1-\sum_i\widehat p_{hi}z_i
                -\sum_{j=0}^4\widehat q_{hj}w_j,
 \qquad
 A_h(z,w)=d_0\left(\sum_i a_{hi}z_i+
                             \sum_{j=0}^4 b_{hj}w_j\right).
\]
At $(e_i,e_j)$, the denominator equals
$d_0(D_{hj}-p_{hi})$, so the score is unchanged. The explicit bounds give
\[
 0<\rho/2<\widehat p_{hi}<2\rho,\qquad
 8\epsilon<\widehat q_{hj}<12\epsilon.
\]
Consequently the sum of all denominator slopes in absolute value is
less than $2r\rho+60\epsilon<61\epsilon<1$. This proves strict
reverse slopes and whole-cube positivity. Letting $\epsilon$ decrease
proves the final compactness assertion. The same $\delta$ works for
all $32$ patterns, so there is no unproved uniformity or numerical
feasibility premise.
\end{proof}

\paragraph{Consequence and remaining gap.}
Every $r\times5$ sign table consequently has an axis-zero strict reverse
compiler with at most nine heads, by the existing Lean-checked theorem
\texttt{hasReverseAxisCompiler\_of\_strictReverseGridCompiler} in
\texttt{Results/ReverseGridAxisRepair.lean}, whose general cost is
$3H+3$. The resulting declarations are
\texttt{hasWeakReverseGridCompiler\_five} and
\texttt{hasReverseAxisCompiler\_five}. Lean checks all rational leading
coefficients and the cleared-numerator identities, and uses continuity
to choose one positive perturbation satisfying all signs and whole-cube
budgets for every row count. The final theorem has no external hypothesis. This does not improve the known
linear-cost compiler asymptotically. Its significance is that the two-head
grid shattering phenomenon persists under arbitrarily strong compactness,
so the finite wide-denominator certificate is not the only source of the
fifth independent bit. A scalable construction with $\Omega(H\log H)$
independent query signs, or another superconstant saving, is still missing.
\end{document}
